\documentclass[11pt]{article}
\usepackage{wbslides}
\deck{Dorigoni: Stokes jumps}
\newcommand{\Sp}{S^{+}\Phi_0}
\newcommand{\Sm}{S^{-}\Phi_0}
\begin{document}

\textslide{Stokes jumps in a two-saddle model}{%
\textbf{Source.} D.~Dorigoni, \emph{An Introduction to Resurgence, Trans-Series and Alien Calculus},
\href{https://arxiv.org/abs/1411.3585}{arXiv:1411.3585v2}, §7, eqs.~166--178, 181, 198.\par\medskip
\textbf{Target.} The two lateral resummations of the perturbative saddle differ by $e^{-1/(2\lambda)}$ times the resummation of the other saddle. Splitting that jump in half on each side cancels the ambiguity and gives a real answer.\par\medskip
\textbf{Conventions.} $\lambda>0$, $q\in\mathbb Z_{\ge0}$, $a=\tfrac12+q$, $b=\tfrac12-q$, $F(z)={}_2F_1(a,b;1;z)$. $S^{\pm}$ = the paper's lateral resummations. $S f = \frac1\lambda\int_0^\infty e^{-t/\lambda}\hat f(t)\,dt$.\par\medskip
\textbf{Imported, not derived.} The lateral contour deformation and the hypergeometric discontinuity (eqs.~172--176); Euler's reflection formula; the thimble decomposition.\par\medskip
\textbf{How to use.} Copy each numbered line and its tag onto your sheet. On a slide with a {\large$\boxed{?}$}, say the line aloud, advance, check, then copy.}

\round{Write the jump as one integral}{Turn the lateral difference into an integral that looks like the other saddle's.}{%
$F(z)={}_2F_1(a,b;1;z)$\\ $a+b=1$\\[0.5ex]
$A=\tfrac{2\pi i}{\Gamma(a)\Gamma(b)}$\\[0.5ex]
$\operatorname{Disc}\hat\Phi_0(t)=A\sqrt{2\pi}\,F(1-2t)$, $t>\tfrac12$\\[0.5ex]
$S\Phi_1=\tfrac{\sqrt{2\pi}}{\lambda}\displaystyle\int_0^\infty e^{-u/\lambda}F(-2u)\,du$}

\given{\Sp-\Sm=\frac1\lambda\int_{1/2}^{\infty}e^{-t/\lambda}\operatorname{Disc}\hat\Phi_0(t)\,dt}{lateral contours}
 {Deform the upper contour onto the lower one. What is left is the integral of the jump across the cut, which starts at $t=\tfrac12$.}
 {This is the paper's exact analytic input. Everything after it is algebra.}

\step{Substitute the discontinuity.}
 {=\frac{A\sqrt{2\pi}}{\lambda}\int_{1/2}^{\infty}e^{-t/\lambda}F(1-2t)\,dt}{cut formula}
 {Replace $\operatorname{Disc}\hat\Phi_0(t)$ by $A\sqrt{2\pi}F(1-2t)$ from the margin. Pull the constants out.}
 {It exposes a hypergeometric function, which we can compare with the other saddle's.}

\choosestep{The cut starts at $\tfrac12$, but $S\Phi_1$ integrates from $0$. Change variable.}
 {=\frac{A\sqrt{2\pi}}{\lambda}\int_{0}^{\infty}e^{-(u+1/2)/\lambda}F(-2u)\,du}{$t=u+\tfrac12$}
 {Limits: $\tfrac12\mapsto0$. Argument: $1-2t=-2u$. Measure: $dt=du$.}
 {The integrand now has the shape of $S\Phi_1$ in the margin.}

\round{Extract the action and the multiplier}{Show the jump is $A\,e^{-1/(2\lambda)}S\Phi_1$ and compute $A$.}{%
$S\Phi_1=\tfrac{\sqrt{2\pi}}{\lambda}\displaystyle\int_0^\infty e^{-u/\lambda}F(-2u)\,du$\\[0.5ex]
$A=\tfrac{2\pi i}{\Gamma(a)\Gamma(b)}$\\ $a=\tfrac12+q$, $b=\tfrac12-q$}

\step{Split the exponential.}
 {=A\,e^{-1/(2\lambda)}\,\frac{\sqrt{2\pi}}{\lambda}\int_0^\infty e^{-u/\lambda}F(-2u)\,du}{$e^{x+y}=e^xe^y$}
 {$e^{-(u+1/2)/\lambda}=e^{-1/(2\lambda)}e^{-u/\lambda}$, and the first factor does not depend on $u$.}
 {It isolates the exponentially small factor: the action of the second saddle.}

\step{Name the remaining integral.}
 {\Sp-\Sm=A\,e^{-1/(2\lambda)}\,S\Phi_1}{def.\ of $S\Phi_1$}
 {What follows $e^{-1/(2\lambda)}$ is exactly $S\Phi_1$ from the margin.}
 {This is eq.~178: the jump is the other saddle, weighted by a multiplier $A$.}

\step{Use $a+b=1$ in $A$.}
 {A=\frac{2\pi i}{\Gamma(a)\Gamma(1-a)}}{$b=1-a$}
 {$a+b=\tfrac12+q+\tfrac12-q=1$.}
 {It puts $A$ in the shape of the reflection formula.}

\step{Apply the reflection formula.}
 {A=2i\sin(\pi a)}{reflection}
 {$\Gamma(a)\Gamma(1-a)=\pi/\sin(\pi a)$, so $2\pi i/(\pi/\sin\pi a)=2i\sin\pi a$.}
 {The Stokes constant no longer has gamma functions.}

\step{Evaluate at $a=\tfrac12+q$.}
 {A=2i(-1)^q}{$\sin(\tfrac\pi2+\pi q)$}
 {$\sin(\tfrac\pi2+\pi q)=\cos(\pi q)=(-1)^q$.}
 {This is the model's coefficient, eq.~181.}

\round{Cancel the ambiguity}{Correct each side by half the jump, so both sides agree.}{%
$E=e^{-1/(2\lambda)}$, $T=S\Phi_1$\\[0.5ex]
$\Sp-\Sm=A\,E\,T$ \ (line 5)\\[0.5ex]
$A=2i(-1)^q$ \ (line 8)\\[0.5ex]
$\Phi_1$ has no jump in this direction (eq.~181).}

\choosestep{Remove the jump symmetrically. What do you subtract on each side?}
 {M^{\pm}=S^{\pm}\Phi_0\mp i(-1)^q\,E\,T}{half the jump}
 {Half of $A$ is $i(-1)^q$. Subtract it above the cut, add it below.}
 {Neither side is preferred, so the corrected value can be real.}

\step{Subtract the two corrected values.}
 {M^+-M^-=(\Sp-\Sm)-2i(-1)^q\,E\,T}{subtract}
 {$\mp$ on each side gives $-\tfrac A2ET-\tfrac A2ET=-AET$.}
 {It puts the jump and the correction in one line.}

\step{Insert the jump.}
 {M^+-M^-=0}{line 5}
 {Line 5 says $\Sp-\Sm=AET=2i(-1)^qET$.}
 {The corrected value does not depend on the side: no ambiguity.}

\step{Read off the saddle parameter in $\Phi_0+\sigma_1E\,\Phi_1$.}
 {\sigma_1^{\pm}=\mp\,i(-1)^q}{match $ET$}
 {Compare $M^\pm$ with $S^\pm(\Phi_0+\sigma_1E\Phi_1)$, coefficient of $ET$.}
 {These are the two choices in eq.~198.}

\round{Reality}{Show the corrected value is real.}{%
$M^{\pm}=S^{\pm}\Phi_0\mp i(-1)^qET$\\[0.5ex]
$E,T$ real for $\lambda>0$\\[0.5ex]
$M^+=M^-$ \ (line 11)}

\given{\overline{\Sp}=\Sm}{conjugation}
 {For real $\lambda$, conjugation swaps the two lateral sums (paper's convention).}
 {}

\step{Conjugate $M^+$.}
 {\overline{M^+}=\Sm+i(-1)^q\,E\,T=M^-}{$\bar i=-i$}
 {}{}

\step{Conclude.}
 {\overline{M^+}=M^+}{line 11}
 {}{The corrected sum is real.}

\textslide{What this did and did not show}{%
\textbf{Shown.} The lateral jump is $2i(-1)^q e^{-1/(2\lambda)}S\Phi_1$ (lines 1--8). Half of it on each side cancels the ambiguity (lines 9--11) and gives a real value (lines 13--15).\par\medskip
\textbf{Not shown.} That this corrected value equals the original integral. That needs the thimble intersection numbers, which are imported.}

\fullchain[8]
\end{document}
